\chapter{Tutorials: the Shell Engine (msg-shell)}
\label{ch:shell}

The \texttt{msg-shell} engine solves structure genes whose heterogeneity lives
in a \emph{surface}: the mid-surface of a thin composite wall. Thickness never
appears in the mesh. It lives inside the wall's Reissner--Mindlin plate law and
inside the through-thickness recovery ladder, so changing a ply schedule is a
change of material data, not a remesh.

Five capabilities ship as runnable examples, all under
\texttt{examples/OpenSG\_shell/}. Every one of them takes a YAML file in and
writes a timed output file in the SwiftComp \texttt{.K} layout.

\begin{table}[htbp]
\centering
\small
\begin{tabular}{llll}
\toprule
example folder & driver & entry point & produces \\
\midrule
\texttt{1\_get\_beam\_props\_...}   & \texttt{beam\_homo\_shell.py}  & \texttt{build\_rm\_bundle}          & Timoshenko $6\times6$ \\
\texttt{2\_get\_beam\_dehom\_...}   & \texttt{beam\_dehom\_shell.py} & \texttt{build\_rm\_bundle} + recovery & 3-D wall stress \\
\texttt{3\_get\_solid\_props\_...}  & four sub-folders               & \texttt{build\_solid\_bundle}       & equivalent solid $6\times6$ \\
\texttt{4\_get\_solid\_props\_...}  & \texttt{schwarz\_solid\_props.py} & \texttt{shell\_sg3d}             & equivalent solid $6\times6$ \\
\texttt{5\_taper\_segment\_...}     & \texttt{run\_taper\_segment.py} & \texttt{segment\_timo\_from\_3dyaml} & tapered-segment $6\times6$ \\
\bottomrule
\end{tabular}
\caption{The five \texttt{msg-shell} capabilities and the example folder that
demonstrates each one. Folder names are abbreviated; the full names appear at
the head of each section below.}
\label{tab:shell-capabilities}
\end{table}

%=====================================================================
\section{Beam properties from a shell cross-section}
\label{sec:shell-beam}

\subsection{What this capability computes}

Given the contour of a thin-walled composite cross-section, this run produces
the $6\times6$ Timoshenko beam stiffness matrix of that section: the constitutive
law a beam solver such as BeamDyn or GEBT needs. It also writes the
$8\times8$ Reissner--Mindlin plate law of every distinct wall region, because
those are the intermediate quantities the beam matrix is assembled from.

The example is \texttt{examples/OpenSG\_shell/1\_get\_beam\_props\_from\_shell\_cross\_section},
and the input is \texttt{iea\_s10\_shell.yaml}: station $r/R = 0.20$ of the
IEA-22 MW reference blade.

\subsection{What a 1-D shell SG is}

A \textbf{1-D shell structure gene} is the contour of the cross-section: a chain
of straight two-node line elements along the wall mid-surface, each element
carrying a laminate stack as material data rather than as geometry. A 2-D solid
cross-section mesh, by contrast, discretizes the wall \emph{through} its
thickness with area elements, so every ply interface must be resolved by the
mesh.

The consequence for the user is size. This station spans about 7.2\,m in $y_2$
and 2.5\,m in $y_3$, has spar caps 87.6\,mm thick and three shear webs, and is
fully described by 307 nodes and 310 line elements. The price is that the walls
are modelled as Reissner--Mindlin plates, so wall-thickness effects enter through
the plate law and the recovery ladder rather than through the mesh.

\subsection{The input file the user edits}

\texttt{build\_rm\_bundle} reads exactly one file. These are its blocks, with the
values found in \texttt{iea\_s10\_shell.yaml}.

\begin{table}[htbp]
\centering
\small
\begin{tabular}{lll}
\toprule
block & content & this station \\
\midrule
\texttt{nodes} & \texttt{[y2 y3 0]} contour coordinates, metres & 307 nodes \\
\texttt{elements} & 1-based two-node line connectivity & 310 elements \\
\texttt{sets: element} & named element sets, one per wall region & 6 sets \\
\texttt{sections} & per set: \texttt{elementSet} plus \texttt{layup} as & 6 sections \\
                  & \texttt{[material, thickness, angle]} triples, outer ply first & \\
\texttt{elementOrientations} & per element, the flattened $3\times3$ frame $[e_1\;e_2\;e_3]$ & 310 rows \\
\texttt{materials} & \texttt{name}, \texttt{density}, \texttt{elastic: \{E, G, nu\}} as 3-vectors & 6 materials \\
\texttt{reference} & the reference surface the whole run is referred to & \texttt{center} \\
\bottomrule
\end{tabular}
\caption{The seven YAML blocks read by \texttt{build\_rm\_bundle}.}
\label{tab:shell-yaml-blocks}
\end{table}

The \texttt{reference} field is the single source of truth. \texttt{build\_rm\_bundle}
reads it and maps it to a thickness fraction (\texttt{center} $\to$ 0.5,
\texttt{oml} $\to$ 0.0, \texttt{iml} and \texttt{oml\_flip} $\to$ 1.0). That one
number is then used consistently for the ring laminate reference, the plate-SG
$z$ origin, the emitted ABD cache and the recovery depth conversion. Pass
\texttt{ref=} explicitly only to override it. Because the reference here is the
mid-surface, quantities that depend on the reference axis, in particular $GJ$ and
the bending--extension couplings, are mid-surface values.

\begin{figure}[htbp]
\centering
\includegraphics[width=0.95\textwidth]{\FIGDIR/shell_xsec_mesh.png}
\caption{The 1-D shell SG of IEA-22 station $r/R = 0.20$: 307 nodes and 310 line
elements along the wall mid-surface, coloured by \texttt{sets: element} entry.
This is the figure written by \texttt{beam\_homo\_shell.py} as
\texttt{iea\_s10\_mesh.png}.}
\label{fig:shell-xsec-mesh}
\end{figure}

Each colour in Figure~\ref{fig:shell-xsec-mesh} is one element set, and each set
carries one laminate stack.

\begin{table}[htbp]
\centering
\footnotesize
\begin{tabular}{llrr}
\toprule
layup & plies, outer ply first & total thickness (m) & elements \\
\midrule
\texttt{layup\_0} & gelcoat, glass\_triax, glass\_uniax, glass\_triax & 0.03260055 & 25 \\
\texttt{layup\_1} & gelcoat, glass\_triax, medium\_density\_foam, glass\_triax & 0.07711403 & 161 \\
\texttt{layup\_2} & gelcoat, glass\_triax, carbon\_uniax (0.08106 m), glass\_triax & 0.08756579 & 13 \\
\texttt{layup\_3} & gelcoat, glass\_triax, glass\_uniax, glass\_triax & 0.02808872 & 8 \\
\texttt{layup\_4} & gelcoat, glass\_triax, carbon\_uniax (0.07707 m), glass\_triax & 0.08358040 & 14 \\
\texttt{layup\_5} & glass\_biax, medium\_density\_foam, glass\_biax & 0.04600000 & 89 \\
\bottomrule
\end{tabular}
\caption{The six laminate stacks of station \texttt{iea\_s10}, as read from the
\texttt{sections} block and echoed in the header of each block of
\texttt{iea\_s10\_shell\_ABDG.out}.}
\label{tab:shell-layups}
\end{table}

\texttt{layup\_2} and \texttt{layup\_4} are the carbon spar caps on the upper and
lower surfaces, \texttt{layup\_1} is the foam-cored skin that covers most of the
contour, \texttt{layup\_5} is the three shear webs, \texttt{layup\_0} is the
trailing-edge reinforcement and \texttt{layup\_3} the leading edge. Every ply
angle in this file is $0^\circ$; the biax and triax fabrics are supplied as
smeared orthotropic materials rather than as individual plies, so the material
frame of the recovery coincides with the wall frame here.

\subsection{Orientation frames}

A wrong \texttt{elementOrientations} block is the most common way to get a
plausible-looking but wrong answer, so the example driver calls
\texttt{auto\_emit} explicitly after the homogenization and writes an orientation
figure. Note that \texttt{auto\_emit} is called \emph{by the driver}, not from
inside \texttt{build\_rm\_bundle}: if you write your own script around
\texttt{build\_rm\_bundle}, add the \texttt{auto\_emit} call yourself.

$e_1$ is the beam axis (out of the page), $e_2$ (blue) is the in-plane ply-flow
direction along the contour, and $e_3$ (green) is the wall normal, which sets the
OML $\to$ IML sense used by the recovery.

\begin{figure}[htbp]
\centering
\includegraphics[width=0.95\textwidth]{\FIGDIR/shell_xsec_orient.png}
\caption{Per-element $e_2$ (in-plane ply direction) and $e_3$ (wall normal) on
the contour, written by \texttt{auto\_emit} as \texttt{iea\_s10\_orient.png}.
Check that the normals point consistently into the section on the skin and
consistently to one side on each web before trusting any stiffness number.}
\label{fig:shell-xsec-orient}
\end{figure}

\subsection{The element and the wall law}

The ring is assembled as a one-quad-deep prismatic strip: the contour is extruded
by the mean element length along the beam axis, and the extruded node row is
DOF-mapped back onto the original row, which enforces the prismatic
(span-invariant) condition exactly. All matrices are then divided by that depth,
so the result is per unit length. The strip actually assembled is returned in the
bundle under the key \texttt{strip} as \texttt{(nodes, quads, h)}.

Each node carries \textbf{six} degrees of freedom,
$[w_1\;w_2\;w_3\;\omega_1\;\omega_2\;\omega_3]$. The drilling rotation
$\omega_3$ is an \emph{independent} field rather than an algebraic function of the
displacements. Making it independent is what allows adjacent walls meeting at a
junction to share a consistent rotation, but it also introduces a spurious mode,
which is removed by a weak drilling constraint carried by one element-constant
Lagrange multiplier per element (\texttt{lam\_space="elem"}). Element-constant is
the inf-sup-stable choice; equal-order nodal multipliers over-constrain as the
mesh is refined.

Transverse shear uses assumed-strain (MITC) tying, and the production ring scheme
(\texttt{shear="mitc4\_g23"}, the default of \texttt{build\_rm\_bundle}) ties
\textbf{only} the $\gamma_{23}$ row. The reason is specific to a prismatic strip:
under span invariance $\gamma_{13}$ carries no fluctuation gradient and is
algebraic in the directors, so tying it would alias the drilling mode instead of
curing shear locking. Only $\gamma_{23}$ pairs a differentiated displacement with
a rotation.

The wall constitutive law is therefore the $8\times8$ Reissner--Mindlin plate law
\begin{equation}
\mathbf{ABDG} =
\begin{bmatrix}
\mathbf{A} & \mathbf{B} & \mathbf{0}\\
\mathbf{B} & \mathbf{D} & \mathbf{0}\\
\mathbf{0} & \mathbf{0} & \mathbf{G}
\end{bmatrix},
\qquad
\bar{\varepsilon} =
\begin{bmatrix}
\varepsilon_{11} & \varepsilon_{22} & 2\varepsilon_{12} &
K_{11} & K_{22} & K_{12}+K_{21} & 2\gamma_{13} & 2\gamma_{23}
\end{bmatrix}^{T},
\label{eq:shell-abdg}
\end{equation}
with exactly that row grouping printed in the header of
\texttt{iea\_s10\_shell\_ABDG.out}:
\texttt{rows [eps11 eps22 2eps12 | K11 K22 K12+K21 | 2g13 2g23]}.

The $2\times2$ transverse-shear block $\mathbf{G}$ has a default,
\texttt{g\_source="msg"}: the MSG/VAM second-order-energy projection (the
Yu--2002 least-squares construction, Eq.~61 of Yu, Hodges \& Volovoi,
\emph{Computers \& Structures} \textbf{81}:439--454, 2003), computed by
\texttt{opensg\_solid.rm\_plate\_1D.msg\_rm\_plate.rm\_plate\_msg}. The
alternative, \texttt{g\_source="whitney"}, is the coupling-aware
complementary-energy shear flow. Table~\ref{tab:shell-gblocks} shows why this is
not a detail on a foam-cored blade section.

\begin{table}[htbp]
\centering
\small
\begin{tabular}{lrr}
\toprule
layup & $G_{11}$ (N/m) & $G_{22}$ (N/m) \\
\midrule
\texttt{layup\_0} & $8.2828865\times10^{7}$ & $9.8804474\times10^{7}$ \\
\texttt{layup\_1} & $4.1457954\times10^{6}$ & $4.1451801\times10^{6}$ \\
\texttt{layup\_2} & $2.8674578\times10^{8}$ & $2.1272143\times10^{8}$ \\
\texttt{layup\_3} & $5.8528099\times10^{7}$ & $8.6668255\times10^{7}$ \\
\texttt{layup\_4} & $2.7301565\times10^{8}$ & $2.0317797\times10^{8}$ \\
\texttt{layup\_5} & $2.4877763\times10^{6}$ & $2.4877763\times10^{6}$ \\
\bottomrule
\end{tabular}
\caption{The MSG wall transverse-shear block $\mathbf{G}$ of each layup, rows 7
and 8 of the $8\times8$ blocks in \texttt{iea\_s10\_shell\_ABDG.out}. The
foam-cored skin \texttt{layup\_1} and the foam-cored webs \texttt{layup\_5} are
roughly seventy and a hundred times more shear-compliant than the carbon cap
\texttt{layup\_2}, so the soft core dominates the section's shear compliance and
must not be assigned by a rule of thumb.}
\label{tab:shell-gblocks}
\end{table}

\subsection{Run it}

The driver has a single user-input line:

\begin{lstlisting}[style=opensg,caption={The user-input block of beam\_homo\_shell.py.}]
############### User Input #################################
name = "iea_s10"               # reads <name>_shell.yaml
############################################################
\end{lstlisting}

\begin{lstlisting}[style=shellst]
cd examples/OpenSG_shell/1_get_beam_props_from_shell_cross_section
python beam_homo_shell.py
\end{lstlisting}

\subsection{The one call}

Everything is behind one function.

\begin{lstlisting}[style=opensg,caption={The minimal script: build the bundle, take the 6x6, emit the orientation figure.}]
import numpy as np
from opensg_shell import build_rm_bundle, auto_emit

B  = build_rm_bundle("iea_s10_shell.yaml")   # RM 6-DOF ring + MSG wall G
C6 = np.asarray(B["Timo"])                   # (6, 6) Timoshenko stiffness
print(B["ref"], B["g_source"])               # "center msg"
print(np.diag(C6))
auto_emit("iea_s10_shell.yaml", out_png="iea_s10_orient.png")
\end{lstlisting}

\texttt{build\_rm\_bundle} returns a \textbf{bundle}, not just a matrix, because
the dehomogenization of Section~\ref{sec:shell-dehom} must reuse the identical
discretization. The bundle carries the keys of
Table~\ref{tab:shell-bundle-keys}.

\begin{table}[htbp]
\centering
\small
\begin{tabular}{ll}
\toprule
key & content \\
\midrule
\texttt{Timo} & the $(6,6)$ Timoshenko stiffness \\
\texttt{V0}, \texttt{V1} & zeroth- and first-order warping fields, $(6m,4)$ each, \\
                         & including the drilling $\omega_3$ boundary values \\
\texttt{corners}, \texttt{red\_cells} & the $(m,2)$ contour coordinates and $(n_{el},2)$ connectivity \\
\texttt{rx3}, \texttt{re3}, \texttt{k22} & 3-D ring coordinates, wall normals and hoop curvature per element \\
\texttt{strip} & the prismatic strip actually assembled, as \texttt{(nodes, quads, h)} \\
\texttt{ax}, \texttt{cross} & the beam-axis index and the two cross-section axis indices \\
\texttt{layup\_per\_elem} & the layup name of each element \\
\texttt{layup\_db}, \texttt{material\_db} & the geometry-free plate-SG databases \\
\texttt{frac}, \texttt{ref}, \texttt{g\_source} & the reference fraction, its name, and the wall-$G$ source \\
\bottomrule
\end{tabular}
\caption{Keys of the dictionary returned by \texttt{build\_rm\_bundle}.}
\label{tab:shell-bundle-keys}
\end{table}

\subsection{What comes out}

\begin{table}[htbp]
\centering
\small
\begin{tabular}{lll}
\toprule
file & written by & content \\
\midrule
\texttt{iea\_s10\_shell\_Timo.out} & \texttt{build\_rm\_bundle} & the effective Timoshenko stiffness and \\
                                   & (via \texttt{write\_sc\_K}) & compliance in the SwiftComp \texttt{.K} layout, \\
                                   & & with an \texttt{OpenSG} banner and a \texttt{Time taken} footer \\
\texttt{iea\_s10\_shell\_ABDG.out} & \texttt{build\_rm\_bundle} & one $8\times8$ RM plate stiffness and compliance \\
                                   & (via \texttt{write\_abdg\_out}) & block per section, each headed by its layup \\
\texttt{abd/iea\_s10\_shell\_abd.yaml} & \texttt{emit\_station\_abd} & per-layup cache of the $8\times8$ wall law plus ply \\
                                   & & lists, thickness and mass per area \\
\texttt{iea\_s10\_mesh.png} & the driver & the contour coloured by layup \\
\texttt{iea\_s10\_orient.png} & \texttt{auto\_emit} & the $e_2$/$e_3$ arrows \\
\bottomrule
\end{tabular}
\caption{Output files of \texttt{beam\_homo\_shell.py}.}
\label{tab:shell-beam-outputs}
\end{table}

Two notes on the ABD cache. It is written \textbf{only if it does not already
exist}, so delete \texttt{abd/iea\_s10\_shell\_abd.yaml} after changing the
layups or the change will not take. And it records \texttt{g\_source: whitney}
and \texttt{reference: mid}, because \texttt{emit\_station\_abd} defaults to the
complementary-energy shear flow. It is a cache for downstream station-level tools
such as shell buckling, not a dump of the ring's own MSG wall law, which lives in
\texttt{\_ABDG.out}.

\subsection{The numbers}

This is the stiffness block of \texttt{iea\_s10\_shell\_Timo.out} exactly as
written, in SI units, on the strain set
$[\varepsilon_{11}\;\gamma_{12}\;\gamma_{13}\;\kappa_1\;\kappa_2\;\kappa_3]$
(extension, the two transverse shears, twist, and the two bending curvatures,
which is the VABS ordering):

\begin{Verbatim}[fontsize=\footnotesize,frame=single,framesep=4pt]
 The Effective Timoshenko Stiffness Matrix
 --------------------------------------------
     2.7660599E+010    -1.7230658E-005     4.6650498E-006     5.4020232E-006     1.6686413E+009    -2.0773040E+009
    -1.7230658E-005     7.1861326E+008     5.7520370E+007    -1.0192044E+008     2.3079953E-006     5.2738749E-006
     4.6650498E-006     5.7520370E+007     4.2174428E+008    -1.6217837E+008     8.4816767E-006     2.7931393E-007
     5.4020232E-006    -1.0192044E+008    -1.6217837E+008     2.4375653E+009     1.6449557E-006    -8.1140553E-006
     1.6686413E+009     2.3079953E-006     8.4816767E-006     1.6449557E-006     3.5143887E+010    -5.5571249E+009
    -2.0773040E+009     5.2738749E-006     2.7931393E-007    -8.1140553E-006    -5.5571249E+009     6.8129724E+010
\end{Verbatim}

The diagonal is what the driver prints.

\begin{table}[htbp]
\centering
\small
\begin{tabular}{llll}
\toprule
term & value & units & meaning \\
\midrule
$EA$    & $2.7660599\times10^{10}$ & N               & axial stiffness \\
$GA_2$  & $7.1861326\times10^{8}$  & N               & transverse shear, chordwise \\
$GA_3$  & $4.2174428\times10^{8}$  & N               & transverse shear, flapwise \\
$GJ$    & $2.4375653\times10^{9}$  & N$\cdot$m$^2$   & torsion about the mid-surface reference \\
$EI_2$  & $3.5143887\times10^{10}$ & N$\cdot$m$^2$   & flapwise bending \\
$EI_3$  & $6.8129724\times10^{10}$ & N$\cdot$m$^2$   & edgewise bending \\
\bottomrule
\end{tabular}
\caption{The Timoshenko diagonal of station \texttt{iea\_s10}, read from
\texttt{iea\_s10\_shell\_Timo.out}.}
\label{tab:shell-timo-diag}
\end{table}

The off-diagonals are the physics that a beam table without couplings would throw
away. The extension--bending entries $1.6686\times10^{9}$ and
$-2.0773\times10^{9}$ N$\cdot$m locate the tension centre away from the
mid-surface reference point; $-5.5571\times10^{9}$ N$\cdot$m$^2$ is the
flap--edge bending coupling, which says the section's principal bending axes are
not the $y_2$/$y_3$ axes; and $-1.0192\times10^{8}$ and $-1.6218\times10^{8}$
N$\cdot$m are the shear--twist couplings that place the shear centre.

Entries printed at $10^{-5}$ against diagonals of $10^{10}$ are exact zeros to
about fifteen digits: extension does not couple to transverse shear or to twist
on this section. That is a useful self-check that the drilling constraint and the
rigid-body kernel are behaving.

The footer of the same file reads \texttt{Time taken: 8.02 sec} for this run.
That covers the whole bundle build: reading the YAML, the six per-layup MSG plate
solves, the ring KKT factorization and the layup-database build.

%=====================================================================
\section{Two-step recovery to 3-D wall stress}
\label{sec:shell-dehom}

\subsection{What this capability computes}

Homogenization gave the beam its stiffness. \emph{Dehomogenization} does the
inverse: given the beam's internal forces at this station, it reconstructs the
3-D stress at every point of every wall, which is what a failure criterion
actually needs. The example is
\texttt{examples/OpenSG\_shell/2\_get\_beam\_dehom\_from\_shell\_cross\_section}
and runs on the same \texttt{iea\_s10\_shell.yaml}.

\subsection{The input the user edits}

\begin{lstlisting}[style=opensg,caption={The user-input block of beam\_dehom\_shell.py.}]
############### User Input #################################
name = "iea_s10"                  # reads <name>_shell.yaml
ff_file = "ff51_rmc_reform.dat"   # beam FF per station (# eta F1 F2 F3 M1 M2 M3)
station_row = 10                  # row 10 -> eta = r/R = 0.20 (iea_s10)
n_depth = 9                       # through-thickness recovery points per element
stress_frame = "material"         # "material" (ply axes) | "plate" (wall axes)
############################################################
\end{lstlisting}

\texttt{ff51\_rmc\_reform.dat} is a plain table whose header line reads
\texttt{\# eta F1 F2 F3 M1 M2 M3 (RM CENTER-ref BeamDyn, VABS order)}: one row per
spanwise station, the first column $\eta = r/R$ and the remaining six the beam
force and moment resultants, in the same VABS order and the same SI units (N,
N$\cdot$m) as the $6\times6$. Row 10 is
\begin{equation}
FF = \begin{bmatrix}
4.1927\times10^{4} & 8.2674\times10^{3} & 1.4673\times10^{6} &
1.2191\times10^{6} & -6.0185\times10^{7} & 3.0426\times10^{5}
\end{bmatrix},
\label{eq:shell-ff}
\end{equation}
a 60\,MN$\cdot$m flapwise moment with a 1.47\,MN flapwise shear, which is what
dominates everything that follows. The forces must be referred to the same axis
as the stiffness: this table is center-referenced, matching
\texttt{reference: center} in the YAML.

\subsection{Run it}

\begin{lstlisting}[style=shellst]
cd examples/OpenSG_shell/2_get_beam_dehom_from_shell_cross_section
python beam_dehom_shell.py
\end{lstlisting}

\subsection{The two-step chain}

\textbf{Step 1, along the contour.} The macro beam strain is
$\mathrm{st} = \mathbf{C}_6^{-1}FF$. That strain is fed through
\texttt{\_macro\_fields}, which recombines the bundle's warping fields into the
nodal warping $w = V_0\,\mathrm{st}_m + V_1\,\mathrm{st}_{c1}$ and its axial
derivative $w' = V_0\,\mathrm{st}_{c1} + V_1\,\mathrm{st}_{c2}$, and then through
\texttt{\_rm\_shell\_strain}, which applies the \emph{same} element operators
used to assemble the stiffness. The result is, per element, the six shell strains
$[\varepsilon_{11}\;\varepsilon_{22}\;2\varepsilon_{12}\;K_{11}\;K_{22}\;2K_{12}]$
at mid-arc. Because it reuses \texttt{quad\_ops\_indep}, step 1 is the exact
energy-consistent adjoint of the homogenization rather than a separate
post-processing model.

The recovery also needs the \emph{gradients} of those strains. The span gradient
$\partial/\partial y_1$ comes from the operators directly; the arc gradient
$\partial/\partial y_2$ is a finite difference of nodally averaged values, and
the averaging is \textbf{layup-boundary aware}. Two element values are averaged
at a node only when both elements carry the same layup, so no difference is ever
taken across a section change where the strains legitimately jump.

\textbf{Step 2, through the thickness.} For each of the 6 layups,
\texttt{rm\_plate\_msg} builds a 1-D MSG-RM plate structure gene (the warping
ladder plus the MSG $\mathbf{G}$), and \texttt{msgrm\_strain\_at\_depth}
evaluates the first-order recovery at a depth $z$ measured from the reference
surface, returning the 3-D Voigt strain, the 3-D Voigt stress
$[\sigma_{11}\;\sigma_{22}\;\sigma_{33}\;\sigma_{23}\;\sigma_{13}\;\sigma_{12}]$
and the ply angle at that depth. Both functions are imported from
\texttt{opensg\_solid.rm\_plate\_1D.msg\_rm\_plate}, shared verbatim with the
plate and solid examples.

The driver evaluates \texttt{n\_depth = 9} ply-interior stations per element, at
$\zeta = (k + \tfrac12)/9$ with $\zeta = 0$ at the OML and $\zeta = 1$ at the
IML, giving $310 \times 9 = 2790$ recovery points.

\begin{lstlisting}[style=opensg,caption={The recovery chain in miniature: one element, one depth.}]
import numpy as np
from opensg_shell import build_rm_bundle, _macro_fields, _rm_shell_strain
from opensg_solid.rm_plate_1D.msg_rm_plate import rm_plate_msg, msgrm_strain_at_depth

B  = build_rm_bundle("iea_s10_shell.yaml")
FF = np.loadtxt("ff51_rmc_reform.dat")[10, 1:]
st, st_m, aA, aB = _macro_fields(B, beam_force_vabs=FF)
s6, s2 = _rm_shell_strain(B, 75, 0.5, st_m, aA, aB)   # element 75, mid-arc

ldb = B["layup_db"]; mdb = B["material_db"]; frac = float(B["frac"])
warp = rm_plate_msg(ldb["layup_2"]["thick"], ldb["layup_2"]["angles"],
                    ldb["layup_2"]["mat_names"], mdb, fraction=frac)
h = float(sum(ldb["layup_2"]["thick"]))
Gam, Sig, ply = msgrm_strain_at_depth(warp, (0.056 - frac) * h,
                                      np.asarray(s6, float), frame="material")
print(Sig / 1e6)      # MPa, Voigt [S11 S22 S33 S23 S13 S12]
\end{lstlisting}

\subsection{What comes out}

\begin{table}[htbp]
\centering
\small
\begin{tabular}{ll}
\toprule
file & content \\
\midrule
\texttt{iea\_s10\_dehom\_shell.txt} & 2790 rows, one per recovery point: \texttt{elem}, \texttt{y2}, \texttt{y3}, \\
 & \texttt{z} from the reference surface, \texttt{zeta}, \texttt{ply} angle in degrees, \\
 & then the six stresses in MPa. The header repeats the frame, $\eta$ and the full \texttt{FF}. \\
\texttt{iea\_s10\_dehom\_stress.png} & the section stress cloud, $\sigma_{11}$ and $\sigma_{12}$ \\
\texttt{iea\_s10\_cap\_tt.png} & $\sigma_{11}$ and $\sigma_{12}$ against $\zeta$ along one through-thickness \\
 & line at the arc-centre element of the peak-stress cap \\
\bottomrule
\end{tabular}
\caption{Output files of \texttt{beam\_dehom\_shell.py}.}
\label{tab:shell-dehom-outputs}
\end{table}

The driver reports the peak as

\begin{Verbatim}[fontsize=\small,frame=single,framesep=4pt]
peak |sigma11| = 318.18 MPa at (y2, y3) = (-0.677, 1.311), zeta = 0.06  [material frame]
\end{Verbatim}

\subsection{Reading the result}

\begin{figure}[htbp]
\centering
\includegraphics[width=0.82\textwidth]{\FIGDIR/shell_dehom_stress.png}
\caption{Recovered $\sigma_{11}$ (top) and $\sigma_{12}$ (bottom) over the
cross-section, all 2790 recovery points, in MPa. Points sit slightly off the
contour because each one is displaced from the mid-surface by its own depth $z$
along the wall normal, which is why the cap points reach $y_3 = 1.311$\,m while
the contour stops at $1.273$\,m.}
\label{fig:shell-dehom-stress}
\end{figure}

The axial stress is confined almost entirely to the two spar caps, exactly as a
flap-dominated load state should be.

\begin{table}[htbp]
\centering
\small
\begin{tabular}{lll}
\toprule
quantity & value & where \\
\midrule
minimum $\sigma_{11}$ & $-318.177$ MPa & element 75, \texttt{layup\_2}, $(y_2,y_3) = (-0.677, 1.311)$, $\zeta = 0.056$ \\
maximum $\sigma_{11}$ & $+314.461$ MPa & element 150, \texttt{layup\_4}, $(y_2,y_3) = (-0.145, -1.249)$, $\zeta = 0.056$ \\
maximum $|\sigma_{12}|$ & $5.573$ MPa & element 11, \texttt{layup\_0}, $(y_2,y_3) = (3.895, 0.335)$ \\
\bottomrule
\end{tabular}
\caption{Extremes of \texttt{iea\_s10\_dehom\_shell.txt} under the load state of
Eq.~\eqref{eq:shell-ff}.}
\label{tab:shell-dehom-extremes}
\end{table}

The upper cap is in compression and the lower cap in tension under this
flap-dominated state, at nearly equal magnitude, and both extremes fall at the
outermost recovery station $\zeta = 0.056$, which for these layups lies just
inside the carbon uniaxial ply. That is where a maximum-stress check on the
carbon belongs. The in-plane shear tells a different story: it peaks not in the
caps but in the trailing-edge and leading-edge regions, where the closed cell
carries the torque as shear flow.

\subsection{What this route recovers, and what it does not}

The repository's \texttt{examples/OpenSG\_shell/IEA\_blade\_beam} benchmark folder
compares this same station's recovery against a VABS 2-D solid run of the
identical section (\texttt{iea\_s10.sg.SM}), and records the result in
\texttt{iea\_s10\_report.txt} as rms differences alongside the rms magnitude of
the VABS field itself.

\begin{table}[htbp]
\centering
\small
\begin{tabular}{llrr}
\toprule
component & region & rms difference (MPa) & rms of the VABS field (MPa) \\
\midrule
$\sigma_{11}$ & skin & 9.529 & 96.676 \\
$\sigma_{11}$ & web & 43.073 & 44.440 \\
$\sigma_{12}$ & web & 6.799 & 22.026 \\
$\sigma_{11}$ & spar-cap through-thickness path & 38.355 & 285.457 \\
$\sigma_{11}$ & cap-left junction path & 1.265 & 286.837 \\
\bottomrule
\end{tabular}
\caption{Two-step shell recovery against a VABS 2-D solid run of the same
section, from \texttt{iea\_s10\_report.txt}.}
\label{tab:shell-dehom-vabs}
\end{table}

The in-plane picture is sound where the load is: about 10\,\% rms on the skin,
13\,\% along the cap thickness, and 0.4\,\% along the reference junction path.
The web $\sigma_{11}$ difference is the same size as the field itself, which is
honest to state plainly: the webs carry almost no axial stress at this state, so
the comparison there is a ratio of two small numbers.

Two components are \textbf{not} delivered by this route and should not be used
from the output file:

\begin{itemize}
\item $\sigma_{33}$ is machine-zero in the first-order recovery (rms
  $3.6\times10^{-13}$ MPa against a VABS rms of $5.6\times10^{-2}$ MPa). This is
  visible directly in the \texttt{S33} column of
  \texttt{iea\_s10\_dehom\_shell.txt}, whose entries are of order $10^{-13}$
  (median $1.96\times10^{-13}$, maximum $9.18\times10^{-13}$).
\item $\sigma_{13}$ is recovered at rms $6.3\times10^{-3}$ MPa against a VABS rms
  of $2.4\times10^{-1}$ MPa. This is structural, not a bug: on a prismatic ring
  the $\gamma_{13}$ row is untied by construction, so it carries the flat-wall
  drilling mode and cannot serve as a recovery target. The companion diagnostic
  \texttt{iea\_s10\_q\_consistency.txt} quantifies this. The tied $\gamma_{23}$
  channel closes its equilibrium identity
  $\int \sigma_{23}\,dz = (\mathbf{G}_{\mathrm{msg}}\,s_2)_2$ to 1.05\,\% median
  error, while the untied $\gamma_{13}$ channel does not, giving a $\sigma_{13}$
  target that is invalid by a factor of order 100. A physical $\sigma_{13}$ needs
  the surface-quad route of Sections~\ref{sec:shell-sg3d}
  and~\ref{sec:shell-segment}, where both rows are tied, not a rescale of this
  ring.
\end{itemize}

The same benchmark folder also records that adding the second-order recovery rung
changes $\sigma_{11}$ by an rms of $1.95\times10^{-4}$ MPa against a first-order
rms of 88.51 MPa, so the first-order chain used by \texttt{beam\_dehom\_shell.py}
is the right default for this class of section.

%=====================================================================
\section{Equivalent 3-D solid from a cross-section SG}
\label{sec:shell-solid-xsec}

\subsection{What this capability computes}

The same 1-D shell structure gene can be homogenized to an \textbf{equivalent 3-D
solid} $6\times6$ instead of a beam $6\times6$: the anisotropic continuum
stiffness of a cell that is tiled periodically in the two transverse directions.
This is the MSG thin-walled (MSG-TW) equivalent-continuum problem, and it is what
you want for a cellular solid, a lattice, or any prismatic microstructure whose
walls are thin.

The strain order is the solid Voigt order
\begin{equation}
\bar{\Gamma} = \begin{bmatrix}
\bar\Gamma_{11} & \bar\Gamma_{22} & \bar\Gamma_{33} &
2\bar\Gamma_{23} & 2\bar\Gamma_{13} & 2\bar\Gamma_{12}
\end{bmatrix}^{T},
\label{eq:shell-gbar}
\end{equation}
with axis 1 along the prismatic direction. That string is returned in the bundle
under the key \texttt{order} and exported as \texttt{GBAR\_ORDER}.

Where the beam problem drives the SG with four beam strains, the solid problem
drives it with the six macro strains, so \textbf{periodicity is always on}:
opposite edges and the corners of the 2-D cell are tied
(\texttt{periodic=True} is the default of \texttt{build\_solid\_bundle}). The wall
law is the same Reissner--Mindlin $8\times8$ of Eq.~\eqref{eq:shell-abdg}, so the
walls carry transverse shear that a classical (CLPT/Kirchhoff) wall model
discards. That distinction is the entire point of the benchmarks below.

\subsection{The driver}

\begin{lstlisting}[style=opensg,caption={The one call for equivalent-solid properties from a cross-section SG.}]
from opensg_shell import build_solid_bundle

b = build_solid_bundle("cell_1Dshell.yaml", cell_area=A_cell)
C3D = b["C3D"]            # 6x6 equivalent solid stiffness, Pa
b["cell_area"], b["area_source"]   # the normalizing area and where it came from
b["solve_time"]           # seconds; also written into <yaml>_C3D.out
\end{lstlisting}

\texttt{cell\_area} is the one argument you must think about: it is the measure
$\omega$ the energy is divided by. Pass it explicitly (then
\texttt{area\_source} reads \texttt{"user"}); if you omit it, the convex-hull
area of the contour is used instead and \texttt{area\_source} reads
\texttt{"hull"}. \texttt{build\_solid\_bundle} writes
\texttt{<yaml>\_C3D.out} by default, in the SwiftComp \texttt{.K} layout
(effective stiffness, effective compliance, orthotropic-approximated engineering
constants) with an \texttt{OpenSG} banner and a \texttt{Time taken} footer, plus
\texttt{<yaml>\_ABDG.out} for the wall laws.

\subsection{Validation: Deo \& Yu (2023)}

The reference is A. Deo and W. Yu, \emph{Equivalent 3D properties of thin-walled
composite structures using mechanics of structure genome}, \emph{Mechanics of
Advanced Materials and Structures} \textbf{30}(9), 2023, 1737--1748. Two columns
are quoted from that paper: \textbf{Deo TW}, its own thin-walled model with a
classical wall law, and the reference \textbf{SwiftComp} column, which is the
paper's MSG solid result and the equivalent of an OpenSG 2-D-solid run. All
values below are MPa, and \texttt{ours \%} and \texttt{Deo TW \%} are both
measured against the SwiftComp column. Every table is the literal content of the
\texttt{*\_compare.dat} file written by the driver.

Four sub-folders live under
\texttt{examples/OpenSG\_shell/3\_get\_solid\_props\_from\_shell\_cross\_section}.
Each is a two-command run: a generator that writes the 1-D shell YAML, then the
driver that homogenizes it.

\begin{lstlisting}[style=shellst]
cd examples/OpenSG_shell/3_get_solid_props_from_shell_cross_section/deo_cellular_solid
python make_deo_cellular_yamls.py
python deo_cellular_solid_props.py
\end{lstlisting}

The generator is where the geometry is edited. For the cellular solid it is a
five-line block at the top of \texttt{make\_deo\_cellular\_yamls.py}:
\texttt{l, h, t = 0.10, 0.10, 0.005}, aluminium \texttt{E, nu = 68.9e9, 0.33},
and \texttt{nseg = 10} elements per wall of length $l$. The driver then computes
the cell area itself from the wall angle,
\texttt{area = (2*l*np.cos(th))*(2*(h + l*np.sin(th)))}, and passes it as
\texttt{cell\_area}.

\subsubsection{Cellular solid, $\theta = +15^\circ$ (paper Table 1)}

\begin{table}[htbp]
\centering
\small
\begin{tabular}{lrrrrr}
\toprule
term & OpenSG-TW (RM) & Deo MSG-TW & SwiftComp & ours \% & Deo TW \% \\
\midrule
$C_{11}$ & 4733.84 & 4736.90 & 4678.90 & $+1.17$ & $+1.24$ \\
$C_{12}$ & 1086.01 & 1089.40 & 1105.50 & $-1.76$ & $-1.46$ \\
$C_{13}$ &  380.63 &  381.81 &  386.88 & $-1.61$ & $-1.31$ \\
$C_{22}$ & 2444.24 & 2446.39 & 2488.90 & $-1.79$ & $-1.71$ \\
$C_{23}$ &  846.69 &  847.44 &  860.89 & $-1.65$ & $-1.56$ \\
$C_{33}$ &  306.74 &  306.99 &  311.48 & $-1.52$ & $-1.44$ \\
$\mathbf{C_{44}}$ & \textbf{4.20} & 4.32 & 4.19 & $\mathbf{+0.18}$ & $+3.09$ \\
$C_{55}$ &  562.61 &  564.15 &  573.11 & $-1.83$ & $-1.56$ \\
$C_{66}$ &  993.78 &  997.52 & 1000.10 & $-0.63$ & $-0.26$ \\
\bottomrule
\end{tabular}
\caption{Deo cellular solid at $\theta = +15^\circ$, MPa, from
\texttt{deo\_cellular\_compare.dat}. Aluminium $E = 68.9$ GPa, $\nu = 0.33$;
$l = h = 10$ cm, $t = 5$ mm. Solve 1.5 s.}
\label{tab:deo-cell-p15}
\end{table}

\subsubsection{Re-entrant cellular solid, $\theta = -15^\circ$ (paper Table 2)}

The auxetic lattice, with the negative $C_{13}$ and $C_{23}$ couplings.

\begin{table}[htbp]
\centering
\small
\begin{tabular}{lrrrrr}
\toprule
term & OpenSG-TW (RM) & Deo MSG-TW & SwiftComp & ours \% & Deo TW \% \\
\midrule
$C_{11}$ & 7505.26 & 7507.40 & 7352.90 & $+2.07$ & $+2.10$ \\
$C_{12}$ & 1090.52 & 1094.10 & 1075.20 & $+1.42$ & $+1.76$ \\
$C_{13}$ & $-219.81$ & $-220.53$ & $-213.94$ & $+2.74$ & $+3.08$ \\
$C_{22}$ & 4151.29 & 4154.90 & 4080.00 & $+1.75$ & $+1.84$ \\
$C_{23}$ & $-846.69$ & $-847.45$ & $-821.77$ & $+3.03$ & $+3.12$ \\
$C_{33}$ &  180.61 &  180.75 &  173.48 & $+4.11$ & $+4.19$ \\
$\mathbf{C_{44}}$ & \textbf{2.47} & 2.68 & 2.47 & $\mathbf{+0.17}$ & $+8.38$ \\
$C_{55}$ &  331.26 &  332.17 &  346.01 & $-4.26$ & $-4.00$ \\
$C_{66}$ & 1687.82 & 1694.20 & 1706.40 & $-1.09$ & $-0.71$ \\
\bottomrule
\end{tabular}
\caption{Deo re-entrant cellular solid at $\theta = -15^\circ$, MPa, from
\texttt{deo\_cellular\_compare.dat}. Solve 0.2 s.}
\label{tab:deo-cell-m15}
\end{table}

$C_{44}$ is the in-plane shear channel, the one the paper flags as its largest
error and attributes to the classical wall model dropping the segment transverse
shear, with a first-order-shear wall model named as future work. That is exactly
the RM wall law used here: the $+3.09$\,\% and $+8.38$\,\% errors drop to
$+0.18$\,\% and $+0.17$\,\%. Every other term tracks the paper's own thin-walled
column to a few tenths of a percent, so the two independent MSG-TW
implementations agree where they share physics.

\subsubsection{Hierarchical square (paper Table 3)}

Run from \texttt{deo\_hierarchical\_square}. The generator
\texttt{make\_deo\_square\_yaml.py} carries \texttt{R, r, t = 0.10, 0.05, 0.005}
and \texttt{nseg = 10}, the same aluminium.

\begin{table}[htbp]
\centering
\small
\begin{tabular}{lrrrrr}
\toprule
term & OpenSG-TW (RM) & Deo MSG-TW & SwiftComp & ours \% & Deo TW \% \\
\midrule
$C_{11}$ & 5183.48 & 5186.39 & 5099.03 & $+1.66$ & $+1.71$ \\
$C_{12}$ &   24.22 &   24.83 &   26.75 & $-9.47$ & $-7.18$ \\
$C_{13}$ &   24.22 &   24.83 &   26.75 & $-9.47$ & $-7.18$ \\
$C_{22}$ &   46.23 &   47.13 &   50.46 & $-8.39$ & $-6.60$ \\
$C_{23}$ &   27.15 &   27.94 &   30.59 & $-11.23$ & $-8.66$ \\
$C_{33}$ &   46.23 &   47.13 &   50.45 & $-8.37$ & $-6.58$ \\
$C_{44}$ &    3.20 &    3.22 &    3.39 & $-5.60$ & $-5.01$ \\
$C_{55}$ &  647.56 &  650.00 &  670.92 & $-3.48$ & $-3.12$ \\
$C_{66}$ &  647.56 &  650.00 &  670.93 & $-3.48$ & $-3.12$ \\
\bottomrule
\end{tabular}
\caption{Deo hierarchical square, MPa, from \texttt{deo\_square\_compare.dat}.
$R = 10$ cm, $r = 5$ cm, $t = 5$ mm. Solve 1.8 s.}
\label{tab:deo-hier}
\end{table}

This lattice is \textbf{bending-dominated}: its transverse channels are
$t^3$-scale ($C_{22}\approx 50$ MPa against $C_{11}\approx 5100$ MPa), so the
load path is wall bending, not wall stretching. Two consequences are worth
recording. First, the RM result is mesh-converged and the offset is physics:
refining from 10 to 40 elements per segment moves $C_{23}$ by 1\,\%, and running
the same model with shear-rigid walls (the Kirchhoff limit) reproduces the
paper's classical column to 0.0--0.3\,\%. The gap between the two columns is
therefore the wall transverse-shear compliance, which is real at the $t/l = 0.1$
ligaments; element order is not involved. Second, both thin-walled models miss
the same junction stiffness. The residual against SwiftComp is the finite
rotational stiffness of the joint regions, which no midline model carries. The
classical model's smaller apparent error is a partial cancellation, artificial
shear rigidity offsetting missing joint stiffness, and the RM model removes the
first error, exposing the second.

\subsubsection{Composite square (paper Fig.~13 / Table 4)}

Run from \texttt{deo\_composite\_square}. The cell is $25.4\times25.4$ mm, so the
driver passes \texttt{cell\_area = L2*L2} with \texttt{L2 = 25.4e-3}; horizontal
segments sit at $y_3=\pm5.84$ mm with one vertical segment through the centre,
and every wall is a $[15]_8$ laminate of $t = 1.016$ mm with $E_1 = 141.96$,
$E_2 = E_3 = 9.79$, $G = 6.136$ GPa, $\nu = 0.42$.

\begin{table}[htbp]
\centering
\small
\begin{tabular}{lrrrrr}
\toprule
term & OpenSG-TW (RM) & Deo MSG-TW & SwiftComp & ours \% & Deo TW \% \\
\midrule
$C_{11}$ & 15263.83 & 16133.90 & 15518.11 & $-1.64$ & $+3.97$ \\
$C_{12}$ &   936.74 &  1016.82 &  1008.76 & $-7.14$ & $+0.80$ \\
$C_{13}$ &   468.37 &   468.38 &   459.01 & $+2.04$ & $+2.04$ \\
$C_{15}$ & $+1191.85$ & $-1192.81$ & $-1135.27$ & sign & $+5.07$ \\
$C_{16}$ &  2383.70 &  2589.74 &  2535.23 & $-5.98$ & $+2.15$ \\
$C_{22}$ &   906.11 &   983.56 &   988.88 & $-8.37$ & $-0.54$ \\
$C_{23}$ &     0.00 &     0.00 &    20.44 & $-100$ & $-100$ \\
$C_{26}$ &   292.25 &   317.09 &   311.80 & $-6.27$ & $+1.70$ \\
$C_{33}$ &   453.06 &   453.06 &   456.25 & $-0.70$ & $-0.70$ \\
$C_{35}$ & $+146.13$ & $-146.06$ & $-140.04$ & sign & $+4.30$ \\
$C_{44}$ &     0.94 &     1.08 &     1.16 & $-18.57$ & $-6.90$ \\
$C_{55}$ &   547.31 &   547.32 &   545.43 & $+0.35$ & $+0.35$ \\
$C_{66}$ &  1094.63 &  1188.20 &  1188.30 & $-7.88$ & $-0.01$ \\
\bottomrule
\end{tabular}
\caption{Deo composite square, MPa, from
\texttt{deo\_composite\_square\_compare.dat}. Solve 1.8 s. The \texttt{.dat} file
prints the two sign-flipped rows as $-204.98$ and $-204.35$ percent, which is the
arithmetic of a sign reversal, not a magnitude error.}
\label{tab:deo-comp}
\end{table}

$C_{13}$, $C_{33}$ and $C_{55}$ reproduce the paper's thin-walled column exactly,
and $C_{11}$ improves on it against SwiftComp. Two open items are recorded
honestly. $C_{15}$ and $C_{35}$ match the published magnitudes to 0.1\,\% with
the \emph{opposite sign}, a ply-angle handedness convention still to be
reconciled, and the likely cause of the $-6$ to $-8$\,\% group. And $C_{23}$ is
zero in both thin-walled models, for the reason below.

\subsubsection{A fifth folder: the isotropic square tube}

\texttt{square\_section\_isotropic} is the simplest case in the set and the best
place to start. It reads \texttt{square\_tube\_1Dshell.yaml} and its user-input
block is four lines: \texttt{a, t\_w = 1.0, 0.03} and
\texttt{CELL\_AREA = 4*a*t\_w}, which is the \emph{wall material area}, the same
measure the solid engines use, so the two routes are directly comparable. Run it
with \texttt{python solid\_props\_msg\_shell.py}; it writes
\texttt{square\_solid\_msg\_shell.out} with the $6\times6$ and the nine
engineering constants obtained from \texttt{elastic\_constants(C)}.

\subsection{Why $C_{23}$ vanishes on axis-aligned lattices}

For a wall at angle $\varphi$ the only $\Gamma_e$ entries carrying weight in
\emph{both} the $\bar\Gamma_{22}$ and $\bar\Gamma_{33}$ columns are the membrane
$\varepsilon_{22}$ row ($\cos^2\varphi$, $\sin^2\varphi$) and the
transverse-shear $2\gamma_{23}$ row ($-\sin\varphi\cos\varphi$,
$+\sin\varphi\cos\varphi$), giving a per-length integrand
\begin{equation}
(A_{11}-G_{\mathrm{msg}})\,\sin^2\varphi\cos^2\varphi .
\label{eq:shell-c23}
\end{equation}
On a $0/90$ lattice every wall has $\sin^2\varphi\cos^2\varphi = 0$, so this
$O(t)$ membrane path vanishes identically and the whole of $C_{23}$ is the
$O(t^2)$ \textbf{junction} contribution: the wall-overlap blocks, which have zero
measure on a midline mesh. Inclined-wall cells are unaffected, which is why the
cellular cases above carry $C_{23}$ correctly and only the square lattices expose
the junction term. Optional junction corrections (\texttt{junction="census"},
\texttt{"micro"}, \texttt{"microcell"}) recover it for stretch-dominated cells;
they do \emph{not} apply to bending-dominated lattices such as the hierarchical
square, and they are off by default.

%=====================================================================
\section{Equivalent 3-D solid from a 3-D shell SG}
\label{sec:shell-sg3d}

\subsection{What this capability computes}

A structure gene periodic in \textbf{all three} directions -- a triply-periodic
minimal surface (TPMS) unit cell, a lattice block, any 3-D microstructure whose
material is a thin sheet -- homogenizes to the same equivalent 3-D solid
$6\times6$ of Eq.~\eqref{eq:shell-gbar}. \texttt{shell\_sg3d} solves it when the
microstructure is a thin sheet meshed as a \emph{surface}: the wall carries the
RM $8\times8$ law, so thickness is a parameter rather than a remesh.

It reuses the cross-section $\Gamma_e$ / $\Gamma_h$ operators unchanged (they are
geometry-general) and adds only the 3-D environment: sparse assembly, the
three-direction periodic map, drilling by penalty on the element-constant
residual, and a three-translation kernel. Shell edges shared by more than two
elements are detected as \textbf{junction lines} and reported as
\texttt{n\_junction\_edges}; a smooth TPMS has none.

The example is
\texttt{examples/OpenSG\_shell/4\_get\_solid\_props\_from\_shell\_3D\_SG}.

\begin{figure}[htbp]
\centering
\includegraphics[width=0.62\textwidth]{\FIGDIR/schwarz_p_mesh.png}
\caption{The Schwarz-P TPMS unit cell as a 3-D shell SG: 13\,536 nodes and
26\,360 triangular shell elements on the midsurface, written by
\texttt{schwarz\_solid\_props.py} as \texttt{schwarz\_p\_mesh.png}. The cell is
the unit cube, and the wall thickness is a number in the YAML, not geometry.}
\label{fig:schwarz-mesh}
\end{figure}

\subsection{The input the user edits}

The mesh comes from a gmsh 2.2 triangle file, and
\texttt{make\_schwarz\_yaml.py} converts it. Its user block is four lines:

\begin{lstlisting}[style=opensg,caption={The user block of make\_schwarz\_yaml.py: mesh path, thickness, material.}]
MSH = ("../../../tests/08072026_square_shell_mesh/TPMS_SC_files/"
       "schwarz_p_D2_shell.msh")
T = 0.036457           # Sample_2, volume-consistent
E, nu = 69.0e9, 0.30
G = E/(2*(1+nu))
\end{lstlisting}

The converter computes each element's frame from the triangle geometry
($e_3$ = unit normal, $e_2$ along the first edge, $e_1 = e_2\times e_3$), writes
one \texttt{sections} entry with the single ply \texttt{[alu, T, 0.0]}, and
reports the surface area and the resulting relative density
\texttt{area.sum()*T}.

\subsection{Run it}

\begin{lstlisting}[style=shellst]
cd examples/OpenSG_shell/4_get_solid_props_from_shell_3D_SG
python make_schwarz_yaml.py
python schwarz_solid_props.py
python compare_boundary_modes.py
\end{lstlisting}

\begin{lstlisting}[style=opensg,caption={The two boundary treatments of shell\_sg3d.}]
from opensg_shell.sg_homo import shell_sg3d

r = shell_sg3d("schwarz_p_3Dshell.yaml")                        # boundary="periodic" (default)
r = shell_sg3d("schwarz_p_3Dshell.yaml", boundary="aperiodic")  # non-periodic SGs (explicit)
r["C3D"], r["n_junction_edges"], r["ndof"], r["solve_time"]
\end{lstlisting}

The SG measure $\omega$ defaults to the \textbf{midsurface area}, the 3-D
analogue of $\omega = \text{perimeter}$ for a plane-section shell SG, and can be
overridden with the \texttt{omega=} argument. The \texttt{.out} file is written
per unit-cell volume so its moduli compare directly with a solid \texttt{.K}.

\subsection{What comes out}

\texttt{shell\_sg3d} writes \texttt{schwarz\_p\_3Dshell\_C3D.out}, and the driver
writes \texttt{schwarz\_p\_mesh.png}. The banner line of the \texttt{.out} records
the whole configuration, and the default periodic run at $t = 0.036457$ gives:

\begin{Verbatim}[fontsize=\footnotesize,frame=single,framesep=4pt]
 OpenSG msg-shell equivalent 3D solid (3-D shell SG, 13536 nodes, 26360 elems,
 0 junction edges, periodic in 3 dirs, per unit cell 1)

 The Effective Stiffness Matrix
 --------------------------------------------
     2.2253086E+009     1.5648739E+009     1.5649503E+009  ...
     1.5648739E+009     2.2252770E+009     1.5647791E+009  ...
     1.5649503E+009     1.5647791E+009     2.2250347E+009  ...

 The Engineering Constants (Approximated as Orthotropic)
 ----------------------------------------------------------
  E1  =     9.3296450E+008
  E2  =     9.3313789E+008
  E3  =     9.3284790E+008
  G12 =     1.0409258E+009
  G13 =     1.0410061E+009
  G23 =     1.0409794E+009
  nu12=     4.1277916E-001
  nu13=     4.1304617E-001
  nu23=     4.1288357E-001

 Time taken: 66.43 sec
\end{Verbatim}

Read the numbers as a health check first. The three normal terms agree to five
digits ($2.2253$, $2.2253$, $2.2250$ GPa) and the three shear terms likewise
($1.04098$, $1.04101$, $1.04093$ GPa), so the effective law is cubic to five
digits, as the Schwarz-P symmetry demands. Every symmetry-forbidden coupling sits
at $10^{3}$--$10^{5}$ Pa against diagonals of $10^{9}$ Pa, four to six orders
below. That is the health check on the all-directions periodic tie: if the tie
map were wrong, those entries would not be small. The repository also records a
head-to-head of this same surface against the msg-solid route on the same
structure, where the shell result runs about 3\,\% below on the normal terms and
about 6\,\% below on the shears, with cubic symmetry intact in both; because the
solid column is digit-identical to SwiftComp, those percentages measure the
thin-shell reduction itself at a surface reaching $t/R \approx 0.2$ near the
saddle necks.

\subsection{Boundary treatment: periodic default versus aperiodic}

\texttt{shell\_sg3d} takes a \texttt{boundary} argument, and the default is
\texttt{"periodic"}: opposite faces, edges and corners of the bounding box are
tied through the sparse assembly map, and the three rigid translations are
removed by area-weighted Lagrange rows. This is the exact treatment for a
periodic medium, and it is what to report.

\texttt{boundary="aperiodic"} is the boundary-solution treatment. For a unit cell
the boundary solution is the macro (affine) field itself, so mapping it onto the
boundary nodes prescribes zero \textbf{translational} warping fluctuation
($w_1=w_2=w_3=0$, Dirichlet) on every node of the bounding-box faces, with the
rotational DOFs left natural. That forbids the affine fluctuations which make a
\emph{free} aperiodic SG rank-one, fixes the rigid modes without a Lagrange
border, and is a kinematic upper bound. Clamping the rotations as well would
roughly double the offset.

\texttt{compare\_boundary\_modes.py} runs both on the current YAML, copies each
per-mode \texttt{.out} to \texttt{schwarz\_t<T>\_<mode>.out}, and appends a block
to \texttt{boundary\_compare.dat}. Table~\ref{tab:schwarz-boundary} is that file.

\begin{table}[htbp]
\centering
\small
\begin{tabular}{lrrr@{\hskip 2em}rrr}
\toprule
 & \multicolumn{3}{c}{$t = 0.036457$} & \multicolumn{3}{c}{$t = 0.129330$} \\
\cmidrule(lr){2-4}\cmidrule(lr){5-7}
term & aperiodic & periodic & \%diff & aperiodic & periodic & \%diff \\
\midrule
$C_{11}$ & 2375.89 & 2225.31 & $+6.767$ & 9881.00 & 8916.81 & $+10.813$ \\
$C_{22}$ & 2375.92 & 2225.28 & $+6.770$ & 9881.00 & 8916.85 & $+10.813$ \\
$C_{33}$ & 2376.22 & 2225.03 & $+6.795$ & 9873.34 & 8916.17 & $+10.735$ \\
$C_{12}$ & 1537.47 & 1564.87 & $-1.751$ & 4921.09 & 5352.21 & $-8.055$ \\
$C_{13}$ & 1539.24 & 1564.95 & $-1.643$ & 4931.10 & 5352.71 & $-7.877$ \\
$C_{23}$ & 1538.99 & 1564.78 & $-1.648$ & 4931.00 & 5352.55 & $-7.876$ \\
$C_{44}$ & 1078.28 & 1040.98 & $+3.583$ & 4150.25 & 3961.63 & $+4.761$ \\
$C_{55}$ & 1078.32 & 1041.01 & $+3.585$ & 4150.26 & 3961.63 & $+4.761$ \\
$C_{66}$ & 1078.61 & 1040.93 & $+3.620$ & 4151.27 & 3961.51 & $+4.790$ \\
\midrule
ndof     & \multicolumn{1}{c}{81216} & \multicolumn{1}{c}{79056} & &
           \multicolumn{1}{c}{81216} & \multicolumn{1}{c}{79056} & \\
solve    & \multicolumn{1}{c}{42.3 s} & \multicolumn{1}{c}{66.5 s} & &
           \multicolumn{1}{c}{42.2 s} & \multicolumn{1}{c}{65.4 s} & \\
\bottomrule
\end{tabular}
\caption{Aperiodic versus periodic on the same Schwarz-P YAML at two wall
thicknesses, per unit cell, MPa. The literal content of
\texttt{boundary\_compare.dat}. In both blocks 720 free-face nodes of 13\,536 are
clamped by the aperiodic treatment.}
\label{tab:schwarz-boundary}
\end{table}

The offset is the clamped-face boundary layer of a \textbf{single} cell: the
expected kinematic bound, not an error in either path. It grows with wall
thickness, from $+6.8$\,\% on $C_{11}$ at $t = 0.036457$ to $+10.8$\,\% at
$t = 0.129330$, because a thicker wall makes the clamped skin a larger fraction of
the cell. Aperiodic is about 35\,\% \emph{faster} here, because the Dirichlet
rows leave the factorization and there is no tie map or Lagrange border, despite
the slightly larger nominal DOF count. Report periodic for a periodic medium;
request \texttt{"aperiodic"} only for SGs that genuinely are not periodic.

%=====================================================================
\section{Tapered and aperiodic segments}
\label{sec:shell-segment}

\subsection{What this capability computes}

Everything so far has assumed the beam is prismatic: the classical msg-shell
analysis meshes only the cross-section contour and enforces periodicity along the
beam axis. A \textbf{tapered} segment has no such periodicity, because its two end
cross-sections differ. \texttt{opensg\_shell.sg\_homo.segment\_timo\_from\_3dyaml}
handles this by replacing the axial periodic condition with Dirichlet data taken
from the two ends. The input is no longer a 1-D ring but a genuine 3-D surface
mesh of a slice of the beam.

The pipeline has four stages:

\begin{enumerate}
\item \textbf{Boundary extraction.} The two end cross-sections are found
  \emph{topologically}: a mesh edge used by exactly one quad is a free edge, and
  each connected component of the free-edge graph with at least three nodes is one
  end ring. The axis is the direction between the two most widely separated
  components, and left/right is decided by the axial coordinate. Each ring is
  written out as a standalone 1-D YAML together with the ring-node $\to$
  segment-node map.
\item \textbf{Boundary solves.} Each ring is homogenized on its own with
  \texttt{ring\_indep}, giving its Timoshenko $6\times6$ and the warping fields
  $V_0$ (Euler--Bernoulli) and $V_1$ (Timoshenko), all six DOFs per node
  including the drilling rotation $\omega_3$.
\item \textbf{Mapping.} The ring $V_0$/$V_1$ are transferred onto the segment's
  boundary nodes as Dirichlet data, replacing the periodic boundary condition and
  fixing the segment's rigid-body modes.
\item \textbf{Segment solve.} One LU factorization of the constrained system
  serves both orders,
  \begin{equation}
  \mathbf{D}_{hh}\,V_0 = -\mathbf{D}_{he}, \qquad
  \mathbf{D}_{hh}\,V_1 = b(V_0),
  \label{eq:seg-orders}
  \end{equation}
  with $V_0 = V_0^{\text{ring}}$ and $V_1 = V_1^{\text{ring}}$ on both ends. The
  generalized-Timoshenko finalization divides the energy by the segment length
  $L$, which is measured as the extent of the nodes along the axis.
\end{enumerate}

The example is
\texttt{examples/OpenSG\_shell/5\_taper\_segment\_aperiodic\_boundaries}.

\subsection{Seeing the capability}

\begin{figure}[htbp]
\centering
\includegraphics[width=0.88\textwidth]{\FIGDIR/aperiodic_segment_bar_urc.png}
\caption{The BAR-URC blade segment the example runs, coloured by its per-element
material frames. This is the surface on which the segment problem is solved: a
3-D shell mesh, not a cross-section.}
\label{fig:seg-bar-urc}
\end{figure}

\begin{figure}[htbp]
\centering
\begin{minipage}{0.48\textwidth}
\centering
\includegraphics[width=\textwidth]{\FIGDIR/aperiodic_ring_L.png}
\end{minipage}\hfill
\begin{minipage}{0.48\textwidth}
\centering
\includegraphics[width=\textwidth]{\FIGDIR/aperiodic_ring_R.png}
\end{minipage}
\caption{The two end cross-sections extracted automatically from the free-edge
graph of the segment mesh, left ring and right ring, each written as a standalone
1-D SG and solved on its own. On a tapered segment the two rings differ, and
their two $6\times6$ matrices bracket the segment result.}
\label{fig:seg-rings}
\end{figure}

\begin{figure}[htbp]
\centering
\includegraphics[width=0.72\textwidth]{\FIGDIR/aperiodic_prismatic_segment.png}
\caption{The prismatic circular cylinder used for the acceptance test. The same
free-edge extraction runs on it, and here both ends are identical by
construction, so the segment $6\times6$ must reproduce the ring $6\times6$.}
\label{fig:seg-prismatic}
\end{figure}

\subsection{The input the user edits}

For the validation case, \texttt{make\_cylinder\_segment.py} generates the
surface mesh. Its parameter block is the whole user interface:

\begin{lstlisting}[style=opensg,caption={The parameter block of make\_cylinder\_segment.py.}]
R  = 1.0            # wall mid-surface radius [m]
NC = 160            # circumferential quads
NL = 3              # axial quads over the segment (>=2 so an interior ring exists)
L  = 1.5            # segment length [m]
HR = [0.05, 0.1, 0.2]   # wall/radius ratios: thin, moderate, thick

ISO = {"name": "iso", "density": 1800.0,
       "E": [70e9, 70e9, 70e9], "G": [26.923e9]*3, "nu": [0.3, 0.3, 0.3]}
ANI = {"name": "ani", "density": 1800.0,
       "E": [37e9, 9e9, 9e9], "G": [4e9, 4e9, 4e9], "nu": [0.3, 0.3, 0.3]}
\end{lstlisting}

It writes six files into a \texttt{meshes/} sub-folder, named
\texttt{seg\_<tag>\_hR<ratio>.yaml}: an isotropic single-ply tube and a balanced
$[+45/-45]$ tube at each of the three thickness ratios. The beam axis is $x$, the
cross-section plane is $(y,z)$, the reference surface is the wall mid-surface,
and each element stores the nine-component triad $[e_1\,e_2\,e_3]$ with $e_1$
axial, $e_2$ the hoop tangent and $e_3$ the inward normal.

\subsection{Run it}

\begin{lstlisting}[style=shellst]
cd examples/OpenSG_shell/5_taper_segment_aperiodic_boundaries
python make_cylinder_segment.py
python run_taper_segment.py
python run_bar_urc_segment.py
\end{lstlisting}

The whole pipeline is one call on the 3-D shell YAML. There are no boundary files
to prepare, because the ends are found from the mesh topology. Note the path: the
cylinder meshes live in the \texttt{meshes/} sub-folder, so the argument carries
that prefix.

\begin{lstlisting}[style=opensg,caption={The segment call, with the meshes/ path used by run\_taper\_segment.py.}]
from opensg_shell.sg_homo import segment_timo_from_3dyaml

r = segment_timo_from_3dyaml("meshes/seg_iso_hR0.1.yaml")
r["S6"]             # segment Timoshenko 6x6
r["C6L"], r["C6R"]  # the two boundary-ring 6x6s
r["L"]              # segment length used to normalize the energy
r["solve_time"]     # seconds
\end{lstlisting}

\subsection{What comes out}

Each run writes, alongside the input YAML, the two boundary 1-D YAMLs
(\texttt{boundary\_<tag>\_L.yaml} and \texttt{\_R.yaml}), the
\texttt{<tag>\_boundaries.npz} bundle, three orientation PNGs
(\texttt{orient\_<tag>\_segment.png}, \texttt{orient\_<tag>\_ring\_L.png},
\texttt{orient\_<tag>\_ring\_R.png}), and the timed \texttt{<yaml>\_Timo.out} in
the SwiftComp layout. For the isotropic cylinder that file opens with

\begin{Verbatim}[fontsize=\small,frame=single,framesep=4pt]
 OpenSG msg-shell tapered/aperiodic segment (boundary V0/V1 Dirichlet, L=1.5)
\end{Verbatim}

\noindent
and closes with \texttt{Time taken: 2.51 sec}.

One caveat on the figures. The extraction is called with \texttt{plot="auto"},
which means an orientation PNG is drawn \emph{only if that file is not already on
disk}. If you regenerate or edit a mesh and rerun, delete the three
\texttt{orient\_*.png} files for that tag first, or you will be looking at the
frames of the previous mesh.

\subsection{Acceptance test: the prismatic identity}

For a \textbf{prismatic} segment the taper terms vanish, so the segment
$6\times6$ must equal the boundary ring's own $6\times6$. That identity is the
acceptance test of the whole boundary extraction plus $V_0$/$V_1$ mapping, and
\texttt{run\_taper\_segment.py} measures it on the prismatic circular tube
($R=1$, $t/R=0.1$, $L=1.5$) in both the isotropic and the $\pm45$ layup, writing
\texttt{seg\_ring\_identity.dat}.

\begin{table}[htbp]
\centering
\small
\begin{tabular}{lrrrl}
\toprule
term & segment $S_6$ & boundary ring $C_6$ & \% error & case \\
\midrule
$EA$   & $4.397947\times10^{10}$ & $4.397947\times10^{10}$ & $0.0000$ & \\
$GA_2$ & $8.473281\times10^{9}$  & $8.473258\times10^{9}$  & $0.0003$ & \texttt{iso\_hR0.1} \\
$GA_3$ & $8.473281\times10^{9}$  & $8.473258\times10^{9}$  & $0.0003$ & $L=1.5$ \\
$GJ$   & $1.694078\times10^{10}$ & $1.694078\times10^{10}$ & $0.0000$ & solve 2.5 s \\
$EI_2$ & $2.200273\times10^{10}$ & $2.200269\times10^{10}$ & $0.0002$ & \\
$EI_3$ & $2.200273\times10^{10}$ & $2.200269\times10^{10}$ & $0.0002$ & \\
\midrule
$EA$   & $7.706160\times10^{9}$  & $7.706160\times10^{9}$  & $0.0000$ & \\
$GA_2$ & $3.262832\times10^{9}$  & $3.262807\times10^{9}$  & $0.0007$ & \texttt{aniso\_hR0.1} \\
$GA_3$ & $3.262832\times10^{9}$  & $3.262807\times10^{9}$  & $0.0007$ & $L=1.5$ \\
$GJ$   & $6.526891\times10^{9}$  & $6.526891\times10^{9}$  & $0.0000$ & solve 1.5 s \\
$EI_2$ & $3.855225\times10^{9}$  & $3.855204\times10^{9}$  & $0.0006$ & \\
$EI_3$ & $3.855225\times10^{9}$  & $3.855204\times10^{9}$  & $0.0006$ & \\
\bottomrule
\end{tabular}
\caption{The prismatic identity, diagonal terms, from
\texttt{seg\_ring\_identity.dat}. The worst diagonal error is $0.0003$\,\% (on
$GA$) for the isotropic tube and $0.0007$\,\% for the $\pm45$ layup. The worst
error anywhere in the $6\times6$ is $0.0052$\,\% and it falls on the
$GA_2$--$EI_2$ shear--bending coupling of the anisotropic case, whose ring value
is $-3.684772\times10^{7}$ against a segment value of $-3.684965\times10^{7}$.}
\label{tab:seg-identity}
\end{table}

The \texttt{\% error} matrix in the file reports insignificant ring terms as
zero: entries whose magnitude is below $10^{-6}$ of the largest are divided by
infinity rather than by numerical noise, which is why the file shows exact zeros
where the ring value is $10^{-5}$ against diagonals of $10^{10}$.

\subsection{Tapered demonstration: BAR-URC segment 12}

\texttt{run\_bar\_urc\_segment.py} runs the same pipeline on a real tapered
segment of the 100 m BAR-URC blade,
\texttt{BAR\_URC\_numEl\_52\_segment\_12.yaml}. Here no identity holds. Each
boundary ring has its own $6\times6$, and the segment result is the true
length-averaged Timoshenko stiffness of the tapered piece. The physical check is
that every segment diagonal falls \emph{between} its two ring values, which the
driver reports as the ratio $(S-R)/(L-R)$; a value in $[0,1]$ means the segment
lies between its two ends.

\begin{table}[htbp]
\centering
\small
\begin{tabular}{lrrrr}
\toprule
term & segment $S$ & ring L & ring R & ratio \\
\midrule
$EA$   & $1.487716\times10^{10}$ & $1.498819\times10^{10}$ & $1.483824\times10^{10}$ & 0.260 \\
$GA_2$ & $6.098229\times10^{8}$  & $6.248039\times10^{8}$  & $6.030954\times10^{8}$  & 0.310 \\
$GA_3$ & $1.398945\times10^{8}$  & $1.457255\times10^{8}$  & $1.345727\times10^{8}$  & 0.477 \\
$GJ$   & $3.353538\times10^{8}$  & $3.684346\times10^{8}$  & $3.057355\times10^{8}$  & 0.472 \\
$EI_2$ & $2.996686\times10^{9}$  & $3.244829\times10^{9}$  & $2.765385\times10^{9}$  & 0.482 \\
$EI_3$ & $8.971707\times10^{9}$  & $9.620926\times10^{9}$  & $8.397279\times10^{9}$  & 0.469 \\
\bottomrule
\end{tabular}
\caption{BAR-URC tapered shell segment 12, from \texttt{bar\_urc\_segment12.dat}.
Segment length $L = 3.45293$ m, solve 15.9 s. All six ratios lie in $[0,1]$, so
every segment diagonal is bracketed by its two boundary rings.}
\label{tab:seg-bar-urc}
\end{table}

Read the ratios as a taper diagnostic. The bending and torsion terms sit near
0.47--0.48, close to the arithmetic mean of the two ends, which is what a nearly
linear taper in those channels produces. $EA$ and $GA_2$ sit at 0.26 and 0.31,
noticeably below the mean, which says the axial and chordwise-shear stiffness of
this segment is closer to its thin end than a linear interpolation of the two
ring values would suggest. A station-by-station beam table built by interpolating
ring properties would therefore over-predict those two terms on this segment.

\subsection{Shear-scheme rule for segments}

\texttt{segment\_timo\_from\_3dyaml} defaults to \texttt{shear="full"}
($2\times2$ Gauss), and that is the production setting to keep: MITC tying
aliases the algebraic drilling content on flat-walled and webbed sections. The
tied variants exposed by the element (\texttt{"mitc4\_wonly"}, which ties both
rows at the Dvorkin--Bathe points but keeps every rotation column at its
full-integration value; \texttt{"mitc4\_g23"}, which ties only the $\gamma_{23}$
row; and \texttt{"mitc4\_both"}, naive full tying) are ablation options for
studying the element, not recommended production settings for segments. Note the
contrast with the ring of Section~\ref{sec:shell-beam}, whose production default
\emph{is} \texttt{"mitc4\_g23"}: the ring is a one-quad-deep prismatic strip and
the segment is a genuine surface mesh, and the two geometries do not want the
same treatment.
